% Restored and adapted detailed content from the previous methodology and compiler implementation draft.

% Source: paper/chapters/04-methodology/conceptual_model.tex

\section{Conceptual Model: Music as a Recursive Hierarchy}\label{detail-sec:conceptual-model}

The fundamental premise of this research is that musical composition is inherently structured and hierarchical.
In traditional music theory, a piece is not merely a sequence of notes, but a complex organization of motifs, phrases,
sections, and movements.
However, many modern digital tools and domain-specific languages treat music as a linear stream of events or a flat
sequence of commands.
This research proposes a shift toward a \textit{structural representation}, where musical ideas are treated as modular,
reusable, and composable units.

\subsection{Linear vs. Structural Representation}\label{detail-subsec:linear-vs.-structural-representation}

In linear representation models, commonly found in live coding environments, the focus is on the \textit{performance}.
Time is often modeled as a continuous thread where commands such as \texttt{play} and \texttt{sleep} are executed
sequentially to generate sound.
While effective for improvisation and real-time interaction, this model often sacrifices the ability to express complex
temporal relationships and hierarchical dependencies.
Modifying a single motif in a linear script often requires broad changes to subsequent timing logic, leading to fragile
and non-modular code.

In contrast, the proposed structural representation models music as a \textit{recursive hierarchy}.
Under this model, a composition is built by nesting smaller musical units within larger ones.
A single note is the base case of this hierarchy.
Notes are grouped into motifs (patterns), motifs are sequenced into phrases, and phrases are assigned to parallel
execution contexts (tracks).
This approach allows for a clean separation of concerns: the musical content (the ``what'') is defined independently of
its temporal placement (the ``when'').

\subsection{Music as a Pure Function}\label{detail-subsec:music-as-a-pure-function}

By treating musical structures as first-class entities, the language allows them to be manipulated using concepts
derived from functional programming.
A musical pattern can be viewed as a \textit{pure function}: it encapsulates a specific musical logic that produces a
deterministic sequence of events.
Just as a function in programming can be called multiple times with different arguments, a musical pattern in
Motivo can be instantiated across different tracks, at different times, or with different parameters (such as
pitch transpositions or rhythmic stretching).

This functional mapping provides several advantages:

\noindent \textbf{Modularity.} Musical ideas can be defined once and reused throughout a composition.

\noindent \textbf{Abstraction.} Complex polyphonic structures can be hidden behind simple pattern interfaces.

\noindent \textbf{Determinism.} Because the language lacks a real-time runtime, every musical relationship is
resolved at compile-time, producing a deterministic representation in the output artifact.

% Source: paper/chapters/04-methodology/formal_grammar.tex

\section{Formal Grammar Specification}\label{detail-sec:formal-grammar}

Motivo is defined by a formal grammar that ensures syntactic consistency and enables the generation of an
Abstract Syntax Tree (AST).
This grammar is divided into a lexical specification (the scanner) and a context-free grammar (the parser).

\subsection{Lexical Grammar}\label{detail-subsec:lexical-grammar}

The lexical layer tokenizes the raw source text into a stream of symbols.
A notable feature of the lexical grammar is the inclusion of musical literals as first-class tokens.
For example, a note literal is defined by the regular expression \verb|[A-G][#b]?[0-8]|, allowing the scanner to
directly recognize pitches such as \texttt{F\#4} or \texttt{Bb2}.
Keywords are reserved for musical structural elements (\texttt{track}, \texttt{pattern}, \texttt{voice}), temporal
modifiers (\texttt{play}, \texttt{from}, \texttt{rest}), and standard programming constructs (\texttt{let},
\texttt{for}, \texttt{if}).

\subsection{Syntactic Grammar}\label{detail-subsec:syntactic-grammar}

The syntactic grammar follows an EBNF (Extended Backus-Naur Form) structure.
It is designed to enforce the hierarchical relationship between musical components.
The root of a program consists of a header (defining metadata like tempo) followed by top-level declarations.

\begin{figure}[htbp]
  \centering
  \begin{lstlisting}[basicstyle=\ttfamily\scriptsize,breaklines=true]
<program>    ::= <header> <top-item>*
<top-item>   ::= <let-stmt> | <pattern-def> | <track-decl>
<track-decl> ::= "track" IDENT? ("using" <inst>)? "{" <track-body> "}"
<track-body> ::= (<stmt> | <pattern-def> | <voice-decl>)*
<pattern>    ::= "pattern" IDENT "(" <params> ")" "{" <stmt-list> "}"
<play-stmt>  ::= "play" <play-target> ";"
<play-target>::= <source> (":" <expr>)? ("from" <expr>)?
  \end{lstlisting}
  \caption{Simplified EBNF representation of the core Motivo grammar.}\label{detail-fig:figure}
\end{figure}

The grammar allows for both braced and unbraced statement bodies in control flow structures, providing a balance between
clarity and brevity.
Crucially, the grammar restricts certain constructs based on their context; for instance, while a \texttt{track} can
contain multiple \texttt{voice} blocks, a \texttt{voice} cannot nest another \texttt{track}, enforcing a one-to-one
mapping between tracks and instrument identities.
The expression sub-grammar supports a full range of arithmetic and logical operators, including a ternary conditional
operator, enabling complex musical logic to be expressed inline.

% Source: paper/chapters/04-methodology/musical_mapping.tex

\section{Mapping Musical Primitives to Programming Constructs}\label{detail-sec:musical-mapping}

To bridge the gap between abstract musical theory and formal computer science, the Motivo language maps
fundamental musical
entities to specific programming constructs.
This mapping ensures that the language remains expressive for musicians while retaining the type safety and structural
clarity of a compiled language.

\subsection{Musical Entities as Data Types}\label{detail-subsec:musical-entities-as-data-types}

The language defines several core primitives that serve as the building blocks of a composition:

\noindent \textbf{Notes.} Represented as pitch-class and octave pairs (e.g., \texttt{C4}, \texttt{F\#2}).
These are treated as constant literals with unique numeric identifiers in the underlying MIDI protocol.

\noindent \textbf{Chords.} Modeled as collections of simultaneous notes.
In Motivo, a chord is a single unit of duration, where the temporal span is determined by its constituent notes.

\noindent \textbf{Sequences.} Represented as ordered lists of musical events (notes, chords, or rests).
This structure maps directly to the concept of a musical ``bar'' or ``cell''.

\noindent \textbf{Rests.} Unlike many languages that use a \texttt{sleep} function to pause execution, this
domain-specific language treats a
rest as a durational value.
A rest is a ``silent note'' that occupies space in the timeline, maintaining the structural alignment of the
composition.

\subsection{Compositional Abstractions}\label{detail-subsec:compositional-abstractions}

The hierarchical organization of music is reflected in the language's scoping and container constructs:

\subsubsection{Patterns: The Unit of Reusability}
A \texttt{pattern} is the primary abstraction mechanism in Motivo\@.
Semantically, it is equivalent to a function definition.
It allows a composer to encapsulate a motif and assign it a name.
Patterns can accept parameters, enabling the creation of dynamic musical ideas---such as a drum fill that
varies based on
a boolean flag or a melodic line that accepts a sequence as a template.

\subsubsection{Tracks: Instrumental Identity}
A \texttt{track} serves as the top-level execution context for a specific instrument.
In the compilation process, each track maps to a unique MIDI channel.
This means that independent musical lines (e.g., a piano melody and a bassline) are kept separate and assigned to
distinct sound patches in the final artifact.

\subsubsection{Voices: Horizontal Parallelism}
Within a single track, the \texttt{voice} construct allows for polyphonic composition.
While a track represents a single ``performer'', a voice represents a single ``line of counterpoint''.
Multiple voices can be defined to start at the same time or at different offsets within the same track, allowing for
complex harmonic structures without the need for manual delta-time calculations.
Voices enable the composer to think in terms of horizontal layers rather than just vertical chords.

% Source: paper/chapters/04-methodology/syntax_design.tex

\section{Syntax Design: C-Style Imperative Flow}\label{detail-sec:syntax-design}

The choice of syntax in a domain-specific language is a critical decision that influences both the learning curve and
the expressiveness of the language.
For Motivo, a C-style imperative syntax was selected for its familiarity and its ability to clearly represent
the deterministic progression of a musical program.

\subsection{Familiarity and Cognitive Load}\label{detail-subsec:familiarity-and-cognitive-load}

By adopting standard programming conventions---such as curly braces for blocks, semicolons for statement
termination, and
parentheses for parameters---Motivo targets a user base that is already familiar with general-purpose languages like C,
C++, or Java.
This choice reduces the cognitive load on the user, allowing them to focus on the musical logic rather than learning a
completely novel syntactical paradigm.
This contrasts with languages like \textit{Tidal}, which rely on deeply nested functional structures that can be opaque
to users without a strong Haskell background.

\subsection{Imperative Flow as a Musical Metaphor}\label{detail-subsec:imperative-flow-as-a-musical-metaphor}

Musical composition is often an imperative process: ``play this, then repeat that, then if a condition is met, change
this''.
Motivo's control flow structures map directly to these compositional needs:

\noindent \textbf{Loops (\texttt{loop}, \texttt{for}):} The \texttt{loop} construct provides a simple mechanism for
repetition (e.g., a constant drum beat).
The C-style \texttt{for} loop enables more complex \textit{algorithmic composition}, where an index variable can
drive musical parameters.
For example, a loop can transpose a motif incrementally by using the iterator value as a pitch offset.

\noindent \textbf{Conditionals (\texttt{if}, \texttt{else}):} These structures allow for \textit{context-aware
composition}.
A phrase might have a different ending depending on whether it is the first or second time it is played, or a track
might include a crash cymbal accent only on the first beat of a section.

\subsection{Algorithmic Variation vs. Static
Definition}\label{detail-subsec:algorithmic-variation-vs.-static-definition}

Traditional MIDI editing is inherently static; a motif is a series of fixed events.
The imperative syntax of Motivo allows motifs to be defined \textit{algorithmically}.
By combining control flow with variables (\texttt{let}), a composer can write code that generates variations of a theme
without manually copying and pasting events.
This methodology reinforces the idea that music is not just a sequence of data points, but a series of logical
instructions that can be computed.
Since the compiler unrolls these instructions at compile-time, the composer enjoys the power of algorithmic generation
without the overhead or latency of a real-time interpreter.

% Source: paper/chapters/04-methodology/temporal_logic.tex

\section{The Timeline and Cursor Model}\label{detail-sec:temporal-logic}

A central challenge in music programming is the representation of time.
While physical time is continuous, musical time is often perceived as a series of discrete units (beats, bars,
subdivisions).
Motivo uses a deterministic \textit{cursor model} to manage temporal placement, providing a balance
between sequential flow and absolute positioning.

\subsection{The Virtual Playhead}\label{detail-subsec:the-virtual-playhead}

The core of the temporal logic is a virtual ``playhead'' or \textit{cursor}.
As the compiler traverses the musical statements within a track or voice, it maintains an internal state representing
the current absolute beat position.
Every musical command interacts with this cursor in one of two ways:

\noindent \textbf{Relative Advancement.} Standard \texttt{play} and \texttt{rest} statements are relative to the current
cursor position.
When a note is played, it is placed at the current cursor value, and the cursor is subsequently advanced by the
duration of that note.
This model allows for natural, sequential composition that mimics the act of writing on a staff.

\noindent \textbf{Absolute Offsets.} The \texttt{from} keyword allows for non-destructive positioning.
A statement such as \texttt{play A4 from 16;} places a note at the 16th beat regardless of the current cursor
position.
Crucially, this operation does not advance the main cursor, allowing the composer to ``jump'' back and forth across
the timeline to layer multiple events without losing the primary sequential flow.

\subsection{Parallelism and Temporal Isolation}\label{detail-subsec:parallelism-and-temporal-isolation}

The cursor model is also applied to parallel structures.
When a \texttt{voice} is defined within a \texttt{track}, it inherits a copy of the parent track's current cursor
position.
However, advancements within the voice are isolated; they do not affect the cursor of the parent track or of sibling
voices.
This allows for the creation of independent rhythmic layers that are synchronized at their start point but proceed at
their own pace.

\subsection{Deterministic Time Resolution}\label{detail-subsec:deterministic-time-resolution}

By using floating-point arithmetic for the internal cursor, the language supports arbitrary rhythmic subdivisions
(e.g., triplets, dotted notes, or complex swing rhythms).
Unlike interpreted systems that may suffer from timing drift or ``jitter'' due to CPU scheduling, the compiled nature of
Motivo calculates every event before writing the output file, so event placement is resolved before playback.
Time is resolved entirely during the ``Lowering'' phase of the compiler, transforming the hierarchical code into a flat,
sorted timeline of events.
This methodology supports that the musical structure is preserved by the current mapping and mathematically consistent.

% Source: paper/chapters/04-methodology/compilation_strategy.tex

\section{The Compilation Strategy: Deterministic Unrolling}\label{detail-sec:compilation-strategy}

Motivo is designed as a fully compiled language, distinguishing it from interpreted musical environments.
This design choice is driven by the mandate for a \textit{zero-runtime architecture}, where all musical and logical
structures are resolved before any sound is produced or any output file is written.

\subsection{The Zero-Runtime Mandate}\label{detail-subsec:the-zero-runtime-mandate}

In a zero-runtime system, the complexity of the composition is handled entirely by the compiler.
This means that features such as nested loops, conditional branching, and recursive pattern calls do not exist in the
output artifact.
Instead, the compiler performs a process of \textit{deterministic unrolling}.
Every loop is expanded into its constituent iterations, every conditional branch is evaluated and resolved, and every
pattern call is inlined.

The result of this process is an Intermediate Representation (IR) that consists of a flat, linear list of atomic musical
events.
This architecture provides several benefits:

\noindent \textbf{High Performance.} The output artifact is extremely lightweight, as all logical overhead has been
removed.

\noindent \textbf{Offline Validation.} Semantic errors---such as variable shadowing, type mismatches, or
invalid rhythmic
math---are caught during the compilation phase, providing immediate feedback to the composer.

\noindent \textbf{Portability.} The resulting IR can be easily serialized into multiple formats (such as MIDI or JSON)
without requiring a specialized interpreter to read it.

\subsection{Temporal Reconstruction of Patterns}\label{detail-subsec:temporal-reconstruction-of-patterns}

A unique aspect of the compilation strategy is how the temporal span of a pattern is reconstructed.
When a pattern is called, the compiler does not just emit the notes within it; it calculates the total duration occupied
by the pattern, including internal gaps and trailing rests.
This is achieved by capturing the state of the internal cursor before and after the pattern's body is unrolled.
The difference between these two points defines the pattern's ``footprint'' on the timeline, allowing the pattern to be
treated as a single, opaque musical unit by the caller.

\subsection{Static Verification of Compositional Logic}\label{detail-subsec:static-verification-of-compositional-logic}

Because the language is compiled, it can enforce constraints that would be difficult to manage in an interpreted system.
The semantic analysis pass verifies that the composition adheres to the physical limits of the target protocol (e.g.,
ensuring notes are within the MIDI pitch range of 0--127) and that the structural hierarchy is valid (e.g., preventing a
\texttt{track} from being defined inside a \texttt{pattern}).
This rigorous validation ensures that the final MIDI artifact is not only syntactically correct but also musically
meaningful and interoperable with other professional audio software.

% Source: paper/chapters/04-methodology/compiler_flow.tex

\section{Compiler Pipeline and Data Flow}\label{detail-sec:compiler-flow}

The architecture of Motivo compiler is organized as a multi-pass pipeline.
Each stage of the pipeline transforms the representation of the musical program, moving from abstract textual intent to
a concrete binary artifact.
By decoupling these stages, the compiler ensures that errors are caught early and that the musical structure is fully
resolved before generation.

\subsection{Information Flow and Artifacts}\label{detail-subsec:information-flow-and-artifacts}

The flow of data through the compiler is illustrated in Figure~\ref{detail-fig:compiler-diagram}.
Each box represents a discrete processing stage, and the annotated arrows describe the intermediate artifacts passed
between them.

\begin{figure}[htbp]
  \centering
  \resizebox{\textwidth}{!}{%
    \begin{tikzpicture}[
        node distance=1.2cm and 2.8cm,
        stage/.style={rectangle, draw=black, thick, fill=blue!5, minimum width=3.5cm, minimum height=1cm,
        align=center, font=\small\sffamily},
        artifact/.style={font=\scriptsize\itshape, text=blue!70!black, align=center},
        arrow/.style={-Stealth, thick}
      ]
      % Stages
      \node (source) [stage, fill=gray!5] {Motivo source\\(.motivo)};
      \node (frontend) [stage, right=of source] {Frontend\\(Scanner \& Parser)};
      \node (semantic) [stage, right=of frontend] {Semantic\\Analyzer};
      \node (lowerer) [stage, below=of semantic] {Lowering\\Engine};
      \node (backend) [stage, left=of lowerer] {Backend\\(MIDI Writer)};
      \node (midi) [stage, fill=gray!5, left=of backend] {MIDI File\\(.mid)};

      % Arrows with Artifacts
      \draw [arrow] (source) -- node[above, artifact] {characters} (frontend);
      \draw [arrow] (frontend) -- node[above, artifact] {AST} (semantic);
      \draw [arrow] (semantic) -- node[right, artifact, xshift=2pt] {Annotated\\AST} (lowerer);
      \draw [arrow] (lowerer) -- node[above, artifact] {IR (Linear Events)} (backend);
      \draw [arrow] (backend) -- node[above, artifact] {binary bytes} (midi);

    \end{tikzpicture}
  }
  \caption{The compiler pipeline showing the sequential transformation of musical
  artifacts.}\label{detail-fig:compiler-diagram}
\end{figure}

\subsection{Transformation Stages}\label{detail-subsec:transformation-stages}

The transformation process can be summarized as follows:

\noindent \textbf{Text to Tree (Frontend).} The raw source text is parsed into an \textit{Abstract Syntax Tree} (AST).
This artifact captures the logical hierarchy of the composition (e.g., nested patterns and tracks) but contains no
musical or temporal resolution.

\noindent \textbf{Tree to Annotated Tree (Semantic Analysis).} The compiler performs symbol resolution and
type inference.
The resulting \textit{Annotated AST} links every identifier to its definition and ensures that all operations are
musically valid.

\noindent \textbf{Tree to Timeline (Lowering).} The hierarchical AST is ``unrolled'' into a flat \textit{Intermediate
Representation} (IR).
This is the most significant transformation, where recursive pattern calls and loops are replaced by a sorted list
of concrete \texttt{NoteEvents} with absolute timestamps.

\noindent \textbf{Timeline to Binary (Backend).} Finally, the IR is serialized.
The backend converts the absolute timestamps into the delta-times required by the MIDI standard and writes the
structured binary stream to disk.

% Source: paper/chapters/04-methodology/midi_target.tex

\section{Choice of Output: MIDI as a Semantic Target}\label{detail-sec:midi-target}

The final component of the methodology is the selection of the output format.
While some musical domain-specific languages target raw audio synthesis or real-time event streaming via Open
Sound Control (OSC), this
research utilizes the \textit{MIDI Format 1} protocol as its primary build target.

\subsection{Interoperability and the Producer Workflow}\label{detail-subsec:interoperability-and-the-producer-workflow}

A key goal of this domain-specific language is to serve as a tool for \textit{producers} and
\textit{composers}, not just a self-contained
instrument.
By targeting MIDI, the compiler produces an artifact that can be understood by many Digital Audio
Workstations, MIDI tools, and notation software.
This interoperability allows the ``code-based composition'' to be seamlessly integrated into a traditional production
workflow, where the resulting MIDI can drive high-end virtual instruments, orchestral libraries, or external hardware
synthesizers.

\subsection{Preserving Symbolic Identity}\label{detail-subsec:preserving-symbolic-identity}

MIDI is a \textit{symbolic representation} of music.
Unlike audio, which is a stream of samples, MIDI preserves the identity of notes, velocities, and rhythmic placement.
This makes MIDI an ideal target for a structured domain-specific language because the high-level constructs
of the language (e.g., a ``C4''
note at beat 2) map directly to low-level MIDI messages (\texttt{Note-On, key 72, velocity 100}).
This mapping keeps the composer's symbolic intent visible in the generated event data,
allowing for further offline analysis or modification in a DAW\@.

\subsection{Offline Analysis and Reproducibility}\label{detail-subsec:offline-analysis-and-reproducibility}

Finally, targeting a static file format like MIDI supports \textit{reproducibility}.
In a live coding environment, the same script might sound different on different machines or under different CPU loads.
With the compiled approach, the compiler generates a deterministic binary file.
This file can be archived, version-controlled, and shared so the compiled event data remains fixed for later inspection
as the composer intended, regardless of the playback environment.
This methodology reinforces the thesis that musical composition is a form of structured data that benefits from the same
rigorous processing as software source code.

% Source: paper/chapters/05-implementation/foundations.tex

\section{Technical Foundations and Infrastructure}\label{detail-sec:implementation-foundations}

The implementation of Motivo compiler is a critical engineering effort that balances high-level musical expressiveness
with low-level binary precision.
This section details the selection of the technology stack and the structural organization of the project, providing a
rationale for the architectural decisions that support the thesis's goals.

\subsection{Language Choice: Modern C++ (C++23)}\label{detail-subsec:language-choice:-modern-c++-(c++23)}

The compiler is written in modern C++, specifically targeting the C++23 standard.
While higher-level languages such as Python or JavaScript are commonly used for domain-specific language
development, C++ was selected for
three primary reasons:

\noindent \textbf{Deterministic Performance.} Musical temporal resolution requires high precision.
By using C++, the compiler can perform recursive unrolling and temporal math with minimal overhead, achieving the
fast local compilation times necessary for real-time feedback in the playground.

\noindent \textbf{Robust Type System.} Modern C++ provides a wealth of tools for building type-safe compilers.
Features such as \texttt{std::variant} and \texttt{std::optional} allow for a functional programming style that is
ideally suited for processing tree-like structures (ASTs).

\noindent \textbf{Binary Mastery.} Since the final build target is a low-level MIDI binary file, C++'s native
support for
bitwise operations and big-endian serialization makes it the most natural choice for implementing the backend.

\subsubsection{Key C++23 Features in the Architecture}

The codebase leverages several advanced features to simplify compiler logic:

\noindent \textbf{\texttt{std::variant} as a Sum Type:} Every AST node and IR value is represented as a variant.
This allows the compiler to treat different musical constructs (e.g., a Note vs a Loop) as a single type while
enforcing exhaustive handling during traversal via \texttt{std::visit}.

\noindent \textbf{Value-Oriented Memory Management.} By combining \texttt{std::variant} with
\texttt{std::unique\_ptr}, the
compiler avoids the ``Object Soup'' anti-pattern.
Nodes are value-like, and recursion is handled via smart pointers that support explicit ownership and reduced
leak risk without the
non-deterministic latency of a garbage collector.

\noindent \textbf{Standard Ranges and Views.} Used extensively in the backend and lowering passes to sort musical events
and filter diagnostic logs efficiently.

\subsection{Modular Build Infrastructure (CMake)}\label{detail-subsec:modular-build-infrastructure-(cmake)}

To ensure scalability and maintainability, the project is structured as a modular hierarchy managed by CMake.
The source code is partitioned into discrete components, each with its own responsibilities and boundaries.

\subsubsection{Library Decomposition}

The compiler logic is distributed across five static libraries:

\noindent \textbf{\texttt{motivo\_common}:} The foundational library containing the core musical enums (\texttt{Pitch},
\texttt{Accidental}), the \texttt{Note} structure, and the \texttt{DiagnosticsEngine}.
It is linked by every other module.

\noindent \textbf{\texttt{motivo\_frontend}:} Wraps the Flex-generated scanner and Bison-generated parser.
It depends on \texttt{motivo\_common} for location tracking and note literals.

\noindent \textbf{\texttt{motivo\_semantic}:} Contains the \texttt{Traversal} engine and symbol table logic.
It takes a raw AST and produces an \texttt{AnalysisResult}.

\noindent \textbf{\texttt{motivo\_lowerer}:} The unrolling engine.
It transforms an \texttt{AnalysisResult} into a flat list of \texttt{NoteEvent} objects.

\noindent \textbf{\texttt{motivo\_backend}:} The MIDI serialization module.
It handles the binary translation of the linear event stream.

\subsubsection{Directory Organization and Public API}

The project follows the ``Pitch-Complete'' directory convention:

\noindent \textbf{\texttt{include/motivo/}:} Contains only the public headers (e.g., \texttt{compiler.hpp}).
This defines the interface for external consumers, such as the CLI wrapper or the playground backend.

\noindent \textbf{\texttt{src/motivo/}:} Contains the private implementation details, including the recursive
visitors and
generated Flex/Bison sources.

This separation ensures that the internal complexity of the compiler pipeline is hidden behind a simple, high-level
API\@.
The top-level \texttt{motivo::compile} function serves as the entry point, orchestrating the sequential flow
of characters
to tree, tree to timeline, and timeline to MIDI\@.

% Source: paper/chapters/05-implementation/frontend.tex

\section{The Frontend: From Text to AST}\label{detail-sec:implementation-frontend}

The frontend is the gatekeeper of the compiler, responsible for validating the syntactic structure of the user's input
and building a logical model of the composition.
This phase is implemented using a classic scanner-parser architecture.

\subsection{Lexical Analysis: Identifying Musical
Literals}\label{detail-subsec:lexical-analysis:-identifying-musical-literals}

The scanner, implemented in \texttt{scanner.l} using Flex, transforms the source code's character stream into a sequence
of tokens.
The most technically significant challenge in the lexical pass was the requirement for first-class musical literals.

\subsubsection{Pitch-Class and Note Recognition}

Unlike conventional languages that treat musical data as strings or numbers, Motivo treats a pitch like
\texttt{F\#4} as
a primitive constant.
The scanner utilizes a specialized regular expression to recognize these patterns:
\begin{lstlisting}[caption={Flex regular expression for note literal recognition.},label={detail-lst:lstlisting}]
pitch        [A-G]
accidental   [#b]
octave       [0-8]
note_lit     {pitch}{accidental}?{octave}
\end{lstlisting}

When this pattern is matched, the scanner does not simply return a token; it invokes
\texttt{parse\_note\_literal}, which
decomposes the string into a \texttt{music::Note} structure (containing the pitch enum, accidental enum, and octave
integer).
This means that the ``musicality'' of the code is validated at the very first byte of processing.

\subsubsection{Keyword and Symbol Disambiguation}

The scanner also manages a set of reserved keywords that define the structural identity of the language.
Keywords such as \texttt{track}, \texttt{pattern}, and \texttt{voice} transition the parser into different scoping
states.
Additionally, the scanner implements ``Longest-Match'' logic to ensure that identifiers starting with keyword names
(e.g., a variable named \texttt{loop\_count}) are not incorrectly tokenized as keywords.

\subsection{Syntactic Analysis: Building the Hierarchy}\label{detail-subsec:syntactic-analysis:-building-the-hierarchy}

The parser, implemented in \texttt{parser.y} using Bison, is a LALR parser with one-token lookahead.
Its primary role is to enforce the context-free grammar and construct the Abstract Syntax Tree (AST).

\subsubsection{AST Node Architecture and Sum-Types}

The AST is designed as a set of recursive structures defined in the \texttt{motivo::ast} namespace.
To achieve maximum technical depth and type safety, the implementation utilizes the ``Sum-Type'' pattern via
\texttt{std::variant}.

\paragraph{Declaration and Definition Nodes}
The top-level structure of the AST is represented by definitions that manage scoping and vertical layering:

\noindent \textbf{\texttt{TrackDefinition}:} Holds an optional name, an optional instrument, and a body of
\texttt{TrackItem} nodes.
A \texttt{TrackItem} is a variant that can hold a statement, a local pattern, or a voice.

\noindent \textbf{\texttt{PatternDefinition}:} A node that captures the signature of a musical function, including its
name, parameter list, and a block of statements.

\noindent \textbf{\texttt{VoiceDefinition}:} Represents horizontal parallelism.
It stores an optional \texttt{from} expression and a body of statements and local patterns.

\paragraph{Statement and Expression Polymorphism}
The leaf nodes of the tree are divided into statements (actions) and expressions (values):

\noindent \textbf{\texttt{StatementKind}:} Contains nodes like \texttt{AssignStatement}, \texttt{ForStatement},
\texttt{IfStatement}, \texttt{LetStatement}, \texttt{LoopStatement}, and \texttt{PlayStatement}.

\noindent \textbf{\texttt{ExpressionKind}:} A variant containing literals (\texttt{Int}, \texttt{Float}, \texttt{Bool},
\texttt{Note}, \texttt{Rest}), identifiers for variable lookup, unary/binary operations, and complex musical
structures like \texttt{SequenceExpression} and \texttt{ChordExpression}.

\subsubsection{Recursive Node Ownership}

To maintain the tree hierarchy, recursive nodes (e.g., the condition of an \texttt{if} statement) are wrapped in
\texttt{std::unique\_ptr}.
This design choice was made to ensure that the AST is a true tree structure in memory, while the use of
\texttt{std::variant} at every level ensures that the compiler can traverse the tree using ``Pattern Matching'' in the
subsequent semantic and lowering passes.
This combination of value-oriented variants and pointer-based recursion provides a robust foundation for the multi-pass
architecture.

% Source: paper/chapters/05-implementation/music_logic.tex

\section{MIDI Domain Logic and Encoding}\label{detail-sec:implementation-music-logic}

A primary objective of the implementation is to bridge the gap between high-level musical notation and the low-level
binary encoding required by the MIDI standard.
This section details the mathematical formulas and mapping strategies implemented in the
\texttt{motivo\_music} module to
ensure that the compiler's output is musically accurate.

\subsection{Note to MIDI Pitch Mapping}\label{detail-subsec:note-to-midi-pitch-mapping}

Motivo allows composers to specify pitches using traditional letter names from A to G, accidentals (\#, b), and
octave numbers
(0--8).
The compiler converts these literals into a single numeric byte (0--127) using a deterministic formula.

The implementation defines the \texttt{music::Note} structure as a combination of three fields:

\noindent \textbf{\texttt{Pitch} enum:} Stores the semitone offset within an octave (\texttt{C=0}, \texttt{D=2},
\texttt{E=4}, \texttt{F=5}, \texttt{G=7}, \texttt{A=9}, \texttt{B=11}).

\noindent \textbf{\texttt{Accidental} enum:} Stores a numeric signed offset (\texttt{Flat=-1}, \texttt{Natural=0},
\texttt{Sharp=1}).

\noindent \textbf{\texttt{Octave} integer:} Stores the numeric octave.

The mapping formula used in the \texttt{Note::midi\_number} function is:
\begin{equation}
  P = 12 \times (O + 2) + V_p + V_a\label{detail-eq:equation}
\end{equation}
Where $P$ is the final MIDI pitch, $O$ is the octave, $V_p$ is the pitch-class value, and $V_a$ is the accidental value.
The constant offset of \textbf{+2} is applied to the octave to align with the standard MIDI convention where \texttt{C4}
is assigned the value 72 (a convention popularized by software such as Ableton Live, meaning ``Middle C'' or \texttt{C3}
is 60).
By centralizing this calculation at the library level, the compiler ensures that all musical logic throughout the
pipeline is absolutely consistent.

\subsection{Instrument and Drum Mapping}\label{detail-subsec:instrument-and-drum-mapping}

The language's \texttt{using} keyword allows the composer to specify the instrumental identity of a track.
This intent is mapped directly to the General MIDI (GM) Level 1 specification.

The implementation utilizes two primary mapping strategies:

\noindent \textbf{Melodic Instruments.} Standard instruments are mapped to their GM program numbers.
For example, \texttt{Piano} maps to 1, \texttt{Guitar} to 24, \texttt{Bass} to 32, and \texttt{Violin} to 40.
These mappings are stored in the \texttt{music::Instrument} enum.
During backend serialization, they are emitted as \textit{Program Change} events on standard MIDI channels (0--8,
11--15).

\noindent \textbf{Percussion Logic.} Unlike melodic tracks, the MIDI standard rigidly reserves \textbf{Channel 10}
specifically for unpitched percussion.
If a track is defined as \texttt{using drums}, the backend automatically routes all its events to Channel 10 and
skips emitting a Program Change event.
Furthermore, individual drum literals in Motivo are mapped to specific GM ``Note-On'' pitches as defined by the
standard drum map:

\noindent \texttt{Kick} $\rightarrow$ 36

\noindent \texttt{Snare} $\rightarrow$ 38

\noindent \texttt{Hihat} $\rightarrow$ 42

\noindent \texttt{Crash} $\rightarrow$ 49

\noindent \texttt{Ride} $\rightarrow$ 51

This explicit mapping supports that the compiled MIDI artifact will sound correct regardless of the Digital Audio
Workstation (DAW) or synthesizer used to play it back.

% Source: paper/chapters/05-implementation/semantic.tex

\section{Semantic Analysis and Validation}\label{detail-sec:implementation-semantic}

After the frontend produces a syntactically correct Abstract Syntax Tree, the compiler enters the \textit{Semantic
Analysis} phase.
This stage is responsible for ensuring that the program is logically and musically sound.
It performs three primary tasks: identifier resolution (scoping), type inference, and musical constraint validation.

\subsection{Architectural Pattern: The Recursive
Visitor}\label{detail-subsec:architectural-pattern:-the-recursive-visitor}

The semantic analyzer is implemented in the compiler's semantic traversal class.
Because the AST uses \texttt{std::variant} for statement and expression polymorphism, the analyzer relies on a visitor
pattern implemented with \texttt{std::visit} and a small overloaded-helper template.
This allows the compiler to dispatch logic cleanly based on the concrete type of the AST node without relying on dynamic
casting or virtual method tables.

The traversal process follows a precise lifecycle designed to support the hierarchical nature of Motivo\@:

\noindent \textbf{Global Collection.} The analyzer first iterates over all top-level global items.
It extracts all \texttt{PatternDefinition} nodes and registers them in the global scope.
This ``pre-pass'' enables forward referencing, allowing a track at the beginning of a file to call a pattern defined
at the end.

\noindent \textbf{Track and Voice Entry.} For each \texttt{TrackDefinition}, the analyzer enters a new lexical scope.
It performs a similar pre-pass to collect track-local patterns before traversing the body statements.
The same logic applies when visiting a \texttt{VoiceDefinition}.

\noindent \textbf{Recursive Resolution.} Expressions are evaluated bottom-up.
A binary expression node, for instance, will recursively visit its left and right children, infer their types,
validate their compatibility, and then determine its own resulting type.

\subsection{Symbol Table and Scope Management}\label{detail-subsec:symbol-table-and-scope-management}

Scoping is managed by the \texttt{SymbolTable} and \texttt{ScopeStack} classes.
The \texttt{SymbolTable} maintains the actual symbol data, while the \texttt{ScopeStack} manages the hierarchical
environment during the AST traversal.

\subsubsection{Hierarchical Scoping Rules}

The \texttt{ScopeStack} maintains a stack of \texttt{ScopeId} references pointing to the \texttt{SymbolTable}.
Scopes are pushed and popped as the visitor enters and exits structural nodes (Global, Track, Voice, Pattern, and
Control Flow Blocks).
The implementation enforces strict visibility rules:

\noindent \textbf{Lookup.} When an \texttt{IdentifierExpression} is encountered, the \texttt{ScopeStack} performs a
reverse search from the current scope up to the global scope.
If the identifier is found, its unique \texttt{SymbolId} and inferred \texttt{Type} are retrieved.

\noindent \textbf{Symbol Metadata.} Every registered symbol stores its \texttt{SymbolKind} (Variable,
  Parameter, Pattern,
Track, or Voice), its declared \texttt{Location} in the source file, and a raw pointer back to its declaration node
in the AST\@.

\subsubsection{AST Annotation}

To separate analysis from execution, the semantic pass does not mutate the original AST structure.
Instead, it populates an \texttt{Annotations} object.
Upon successful resolution of an identifier, the \texttt{Annotations} object maps the memory address of the
\texttt{IdentifierExpression} to its resolved \texttt{SymbolId}.
Similarly, every evaluated expression is mapped to its inferred \texttt{Type}.
This metadata ensures that the subsequent ``Lowering'' pass does not need to perform expensive string lookups or type
checks; it simply retrieves the cached resolution data.

\subsection{Formal Type Rules and Constraints}\label{detail-subsec:formal-type-rules-and-constraints}

The type system acts as the ``Musical Guardrail'' of the compiler, implemented in the \texttt{type\_rules.cpp} module.

\subsubsection{Type Categories}
The analyzer categorizes types into the \texttt{TypeKind} enumeration:

\noindent \textbf{Scalars.} \texttt{Int}, \texttt{Double}, \texttt{Bool}.
These are utilized for mathematical operations, loop counts, and conditional logic.

\noindent \textbf{Musical Primitives.} \texttt{Note}, \texttt{Rest}, \texttt{Chord}, \texttt{Sequence}, \texttt{Drum}.
These represent the playable entities that occupy space on the timeline.

\noindent \textbf{Utility Types.} \texttt{Unknown} (used during error recovery) and \texttt{Void} (for statements that
produce no value).

\subsubsection{Constraint Validation Logic}
The analyzer enforces critical invariants through specific validator methods within the \texttt{Traversal} class:

\noindent \textbf{\texttt{validate\_binary\_operands}:} Ensures that arithmetic operators (\texttt{+}, \texttt{-},
\texttt{*}, \texttt{/}, \texttt{\%}) are strictly applied to numeric scalars (\texttt{Int} or \texttt{Double}).
For example, attempting to add two sequences (\texttt{[C4] + [D4]}) will trigger a semantic error, as musical
concatenation is handled differently than arithmetic addition.

\noindent \textbf{\texttt{validate\_play\_target}:} This method verifies that the target of a \texttt{play}
statement resolves to a valid musical primitive.
If a composer mistakenly writes \texttt{play (i < 5);}, the analyzer intercepts the error because the target
evaluates to a \texttt{Bool}.

\noindent \textbf{\texttt{validate\_pattern\_instantiation}:} When a \texttt{PatternCallExpression} is visited,
the analyzer
performs a virtual instantiation.
It creates a temporary scope, binds the parameter names to the types of the arguments provided at the call site, and
recursively visits the pattern's body block.
This allows the compiler to catch type errors \textit{inside} a pattern that only manifest when it is invoked with
specific argument types, ensuring deep semantic validation.

% Source: paper/chapters/05-implementation/lowering.tex

\section{The Lowering Engine: Temporal Resolution}\label{detail-sec:implementation-lowering}

The \textit{Lowering Engine} is the most sophisticated component of the compiler.
It performs the transformation from the hierarchical, logically structured Annotated AST into a flat, chronologically
sorted list of events known as the Intermediate Representation (IR).
This process ``unrolls'' the recursive structure of the music and resolves all temporal logic into absolute beat
offsets, effectively stripping away all programming control flow.

\subsection{The Intermediate Representation (IR)}\label{detail-subsec:the-intermediate-representation-(ir)}

The output of the lowering engine is an \texttt{ir::Program} object.
This artifact represents the physical timeline of the composition.
Its primary building block is the \texttt{NoteEvent} structure:

\begin{lstlisting}[language=C++, caption={Structure of an atomic IR event.},label={detail-lst:lstlisting2}]
struct NoteEvent {
    int midi_note;          // Numeric pitch (0-127)
    double start_beat;      // Absolute position from track start
    double duration_beats;  // Temporal span in beats
    int velocity;           // MIDI intensity (default 100)
};
\end{lstlisting}

A key architectural detail is that the IR is entirely agnostic of the language's original control flow.
Whether a note was generated by a static \texttt{play} command, a \texttt{for} loop iteration, or a recursive
\texttt{pattern} call, it is represented in the IR as a simple, independent data point.

\subsection{Stateful Unrolling and Cursor Management}\label{detail-subsec:stateful-unrolling-and-cursor-management}

The unrolling engine is managed by the \texttt{LowererContext} class.
Unlike the semantic analyzer, which checks static types, the lowerer acts almost as an interpreter, executing the logic
and managing actual runtime values.
It maintains a ``cursor'' (a \texttt{double} passed by reference) tracking the current playhead position in beats.

The \texttt{lower\_statement} function utilizes \texttt{std::visit} to dispatch logic based on the statement kind:

\noindent \textbf{\texttt{PlayStatement}:} The lowerer evaluates the play target.
It determines the note duration (using the default, or overriding it if the \texttt{:} operator is present).
The start beat is set to the current cursor value, and the event is emitted.
Finally, the cursor is advanced by the duration.
If the \texttt{from} keyword is used, the start beat is overridden by the evaluated offset, and the cursor is
\textit{not} advanced.

\noindent \textbf{\texttt{LoopStatement} and \texttt{ForStatement}:} The engine evaluates the iteration count or
condition, then enters a loop in C++, recursively calling \texttt{lower\_block} for the body.
Because the \texttt{cursor} reference is passed down, the temporal state is naturally preserved across iterations,
placing notes end-to-end.

\noindent \textbf{\texttt{IfStatement}:} The lowerer evaluates the boolean condition into a concrete \texttt{true} or
\texttt{false}.
Only the matching branch is traversed, effectively performing ``compile-time pruning''.

\subsection{Dynamic Expression Evaluation}\label{detail-subsec:dynamic-expression-evaluation}

During the lowering phase, expressions are resolved into concrete values represented by the \texttt{ir::Value} variant.
These values can be scalars (\texttt{int}, \texttt{double}, \texttt{bool}) or musical structures (\texttt{NoteValue},
\texttt{RestValue}, \texttt{SequenceValue}, \texttt{ChordValue}).

The \texttt{LowererContext} maintains a \textit{Dynamic Scope Stack}: a vector of hash maps from identifier names to
\texttt{ir::Value} entries.
When a \texttt{let} or assignment statement is encountered, the evaluated \texttt{ir::Value} is bound to the identifier.
This allows for dynamic algorithmic generation, where loop indices or accumulator variables can be used as multipliers
for note durations or offsets.

\subsection{Value Flattening and Complex Structures}\label{detail-subsec:value-flattening-and-complex-structures}

A significant challenge in the lowering engine is the conversion of complex musical types into atomic
\texttt{NoteEvent}s.
This is handled by the \texttt{ValueFlattener} module.

\subsubsection{Flattening Chords}
When a \texttt{ChordValue} is encountered, the flattener iterates over all its constituent notes and emits them using
the \textit{same} \texttt{start\_beat}.
The main lowering routine then advances the track's cursor by the \textit{maximum} duration found among the chord's
members, ensuring the next musical statement begins only after the entire chord has finished sounding.

\subsubsection{Pattern Temporal Reconstruction}
The most technically complex feature of the lowerer is how it preserves the ``temporal integrity'' of patterns.
When a \texttt{PatternCallExpression} is evaluated:

\noindent The lowerer snapshots the current cursor position: \texttt{start\_time = cursor}.

\noindent It pushes a new variable scope, binds the arguments to the parameter names, and unrolls the pattern's body.

\noindent It captures the final cursor state: \texttt{end\_time = cursor}.

\noindent It calculates the total temporal span as the difference between \texttt{end\_time} and \texttt{start\_time}.

To ensure that internal gaps and trailing rests within the pattern are respected, the lowerer rebuilds a
\texttt{SequenceValue} that intersperses the emitted notes with \texttt{RestValue} items matching the exact temporal
gaps.
This reconstructed value is then returned to the caller, allowing the pattern to be treated as a single, opaque rhythmic
block.

\subsection{Voice Parallelism and Cursor Isolation}\label{detail-subsec:voice-parallelism-and-cursor-isolation}

Parallel melodic lines are managed through the \texttt{voice} construct.
To implement parallelism without shared state conflicts or manual timing math, the lowerer uses
\textit{Cursor Isolation}:

\noindent Upon entering a \texttt{VoiceDefinition}, the lowerer captures the current outer cursor value.

\noindent It creates a local \textit{copy} called \texttt{voice\_cursor}.
If the voice specifies a \texttt{from} expression, the \texttt{voice\_cursor} is initialized to that value instead.

\noindent All statements unrolled within the voice modify only the \texttt{voice\_cursor}.

\noindent The generated \texttt{NoteEvent}s are appended to the track's master list, but upon exiting the
voice block, the
\texttt{outer\_cursor} is completely unmodified.

This isolation supports that the parent track and other sibling voices remain aligned according to the
compiler cursor model at their starting
beats, regardless of the complexity or length of the parallel voice.

% Source: paper/chapters/05-implementation/backend.tex

\section{The MIDI Backend: Binary Serialization}\label{detail-sec:implementation-backend}

The final stage of the compilation pipeline translates the flat list of linear \texttt{NoteEvent}s (the IR) into a
Standard MIDI File (SMF) Format 1 binary artifact.
This process requires careful management of binary structures, delta-time calculations, and variable-length encodings.

\subsection{SMF Format 1 Structure}\label{detail-subsec:smf-format-1-structure}

The backend generates a multi-track MIDI file structured into binary chunks.
Each chunk begins with a 4-byte ASCII identifier followed by a 32-bit big-endian integer denoting the chunk's data
length.
The compiler enforces strict big-endian serialization via custom bit-shifting utility functions (\texttt{write\_u16\_be}
and \texttt{write\_u32\_be}).

\noindent \textbf{\texttt{MThd} (Header Chunk):} The file begins with a header specifying \texttt{Format=1}
(multiple simultaneous tracks).
The compiler calculates the total number of tracks as $1 + N$ (where $N$ is the number of instrumental tracks in the
Motivo source) to account for the mandatory tempo track.
The temporal resolution is hardcoded to \textbf{480 PPQ} (Pulses Per Quarter-note), providing a high degree of
precision for complex rhythmic subdivisions.

\noindent \textbf{\texttt{MTrk} (Track Chunks):} The first track (Track 0) is initialized as the ``Tempo Track''.
It contains a \texttt{0xFF 51 03} meta-event.
The backend calculates the \textit{Microseconds Per Quarter-note} ($uspb$) using the formula $uspb =
\frac{60,000,000}{\mathrm{BPM}}$ and writes this as a 24-bit value.
Following the tempo track, each domain-specific language track is serialized into its own \texttt{MTrk} chunk.

\subsection{The Delta-Time Sorting Algorithm}\label{detail-subsec:the-delta-time-sorting-algorithm}

A fundamental constraint of the MIDI protocol is that it does not utilize absolute timestamps.
Instead, every event within a track must be preceded by a \textit{delta-time}, representing the number of ticks to wait
\textit{after} the immediately preceding event.

Because the IR unrolling process can place events into the track out-of-order (for example, when a \texttt{voice} block
places notes concurrently with the main track via the \texttt{from} keyword), the backend must perform a
re-serialization pass:

\noindent \textbf{Tick Conversion.} Every absolute beat position and duration in the IR is multiplied by the
PPQ constant
to convert it into absolute ticks (\texttt{ticks = beat * 480}).
This creates two tick points for every note: an ``on-tick'' and an ``off-tick''.

\noindent \textbf{Event Separation.} The backend transforms each \texttt{NoteEvent} into two separate
low-level events: a
\texttt{Note-On} event at the start tick, and a \texttt{Note-Off} event at the end tick.

\noindent \textbf{Stable Sorting.} All separated events within a track are collected into a vector and sorted using
\texttt{std::ranges::stable\_sort} by their absolute tick value.
The stable sort supports that if a \texttt{Note-Off} and a \texttt{Note-On} for the same pitch occur at the exact
same tick, their relative ordering is preserved, preventing overlapping silence issues.

\noindent \textbf{Delta Computation.} The backend iterates sequentially through the sorted list.
For each event, the delta is computed as the difference between \texttt{current\_tick} and
\texttt{previous\_tick}.

\subsection{Variable-Length Quantities (VLQ)}\label{detail-subsec:variable-length-quantities-(vlq)}

Delta-times in MIDI are compressed using \textit{Variable-Length Quantities}.
This 7-bit encoding scheme allows small delta-times (e.g., zero ticks for simultaneous chord notes) to occupy only one
byte, while larger delays can expand up to four bytes.
The compiler implements a custom \texttt{write\_vlq} function that continuously shifts the 32-bit tick delta by 7 bits,
storing the chunks in a buffer, and applying the high-bit continuation flag (hexadecimal \texttt{80}) to all but the
final byte.

\subsection{Event Emission}\label{detail-subsec:event-emission}

Finally, the backend emits the actual status and data bytes for every event:

\noindent \textbf{Program Change (\texttt{0xC0 | channel}):} Emitted at tick 0 for every melodic track to assign the
chosen instrument.

\noindent \textbf{Note-On (hexadecimal \texttt{90} with the channel encoded in the low nibble):} Contains the pitch and
the defined velocity.

\noindent \textbf{Note-Off (hexadecimal \texttt{80} with the channel encoded in the low nibble):} Explicitly
emitted with
a velocity of 0 to terminate the sound.

\noindent \textbf{End of Track (\texttt{0xFF 2F 00}):} A mandatory meta-event appended to the end of every \texttt{MTrk}
chunk to signal the end of the data stream.

By strictly separating the absolute temporal logic of the IR from the relative delta-time mechanics of the MIDI backend,
the compiler ensures that the musical timeline is mathematically sound and that the resulting binary file strictly
adheres to the standard protocol.
